\section{Conclusiones}
\label{sec:conclusiones}

Este trabajo describió EAM-TMS, un sistema de memoria transactiva multi-agente construido sobre la Memoria Asociativa Entrópica (EAM) y su extensión
hetero-asociativa (EHAM)~\cite{pineda2021entropic,morales2024hetero}.
El sistema opera de forma bidireccional texto $\leftrightarrow$ imagen con rechazo claro en ambas direcciones, sin recurrir a módulos externos de clasificación o umbrales supervisados.

\paragraph{Logros principales.}
El sistema alcanzó una \textbf{precisión madura del 98.75\%} en el banco de 80 consultas semánticas (texto $\to$ imagen) con una tasa de rechazo del 1.25\%, y una \textbf{fidelidad del 97.5\%} entre las fases temprana y madura.
En la dirección inversa (imagen $\to$ palabras), el hemisferio visual logró una \textbf{precisión de enrutamiento del 100\%} sobre las imágenes aceptadas y una \textbf{evocación de dominio del 94.1\%} (16/17 imágenes produjeron al menos una etiqueta semántica correcta).

El directorio transactivo se formó por interacción sin supervisión explícita, alcanzando un balance de \textbf{98.9\%} del máximo teórico (1.567 de 1.585 bits) con conteos $[78, 65, 53]$.
Con presentación intercalada, la transición al 90\% de precisión ocurrió en \textbf{$k = 13$ interacciones}, lo que muestra que el directorio aprende rápido cuando la exposición es diversa.

La especificidad del sistema fue \textbf{perfecta}: en ningún valor de $N$ probado se aceptó una imagen falsa (0\% en todo el rango). Esto indica que el containment inverso con $\xi = 0$ es un criterio de aceptación efectivo y conservador.

\paragraph{Lectura general.}
El punto central del trabajo es que el comportamiento anómalo observado inicialmente (sesgo de densidad, recuperaciones sin variedad) no era un defecto del modelo EAM sino consecuencia de un llenado con prototipos promedio.
Al alimentar el sistema con instancias reales, el modelo de Pineda y Morales opera correctamente sin más ajustes externos.
Esto refuerza la importancia de respetar las condiciones de uso del modelo antes de buscar parches del sistema propuesto.

\paragraph{Trabajo futuro.}
Los siguientes pasos más relevantes son:
\begin{enumerate}
    \item \textbf{Barrido de $\xi > 0$}: incrementar la cobertura del hemisferio visual manteniendo la especificidad perfecta.
    \item \textbf{Magnitudes de fastText}: reemplazar la binarización por signo por una representación ponderada por magnitud para ampliar el margen discriminativo intra/inter-clase (actualmente de 3.3 puntos).
    \item \textbf{Cuarto agente y rotación de miembros}: explorar olvido controlado y reorganización del directorio cuando un agente abandona o se incorpora al sistema.
    \item \textbf{Más dominios por agente}: permitir que un mismo agente aprenda varias clases o dominios, en lugar de mantener una sola categoría visual por agente.
    \item \textbf{Aprendizaje en fase madura}: extender el sistema para que los agentes sigan almacenando labels, imágenes y asociaciones nuevas después de formar el directorio, cubriendo procesos como \emph{transactive encoding} y \emph{transactive retrieval}.
    \item \textbf{Evocación de atributos}: extender el hemisferio visual para que una imagen de una manzana verde evoque la etiqueta ``green'', emparejando labels con regiones de la imagen mediante atributos físicos durante el llenado.
\end{enumerate}

\bigskip

En resumen, los resultados muestran que las Memorias Asociativas Entrópicas de Pineda y Morales son una base suficiente para construir un sistema de memoria transactiva multi-agente funcional, bidireccional y con rechazo controlado, sin recurrir a modelos de aprendizaje profundo supervisado para la toma de decisiones.
